\chapter{Introdução}
\label{cap:introducao}

% Seções primárias: iniciam em nova página, em negrito e maiúsculas
% (a conversão para maiúsculas é automática).

[Apresente o tema, o problema de pesquisa, a justificativa e a estrutura do trabalho.]

Este modelo segue os templates da \gls{UFERSA} e as normas da \gls{ABNT}. O primeiro uso do comando \texttt{\textbackslash gls\{SIGLA\}} imprime o nome por extenso seguido da sigla; os usos seguintes imprimem apenas a sigla, como em \gls{ABNT}. Somente as siglas usadas no texto entram na lista de abreviaturas e siglas.

\section{Objetivos}

% Seções secundárias: negrito, com maiúsculas e minúsculas conforme a
% gramática. Digite o título como deve aparecer.

\subsection{Objetivo geral}

% Seções terciárias a quinárias: sem negrito.

[Descreva o objetivo geral do trabalho.]

\subsection{Objetivos específicos}

Exemplo de alíneas (NBR 6024): o texto que antecede as alíneas termina em dois-pontos, e cada alínea começa com letra minúscula e termina em ponto e vírgula, exceto a última, que termina em ponto:

\begin{alineas}
	\item primeiro objetivo específico;
	\item segundo objetivo específico;
	\item terceiro objetivo específico.
\end{alineas}

\section{Estrutura do trabalho}

O Capítulo~\ref{cap:fundamentacao-teorica} apresenta a fundamentação teórica; o Capítulo~\ref{cap:metodologia}, a metodologia; o Capítulo~\ref{cap:resultados}, os resultados; e o Capítulo~\ref{cap:conclusao}, as considerações finais.
